\documentclass[11pt]{article}

\usepackage[margin=1in]{geometry}
\usepackage[T1]{fontenc}
\usepackage[utf8]{inputenc}
\usepackage{lmodern}
\usepackage{hyperref}
\usepackage{enumitem}
\usepackage{longtable}
\usepackage{array}
\usepackage{booktabs}
\usepackage{xcolor}
\usepackage{titlesec}

\hypersetup{
    colorlinks=true,
    linkcolor=black,
    urlcolor=blue
}

\titleformat{\section}{\Large\bfseries}{\thesection.}{0.5em}{}
\titleformat{\subsection}{\large\bfseries}{\thesubsection}{0.5em}{}
\titleformat{\subsubsection}{\normalsize\bfseries}{\thesubsubsection}{0.5em}{}

\newcommand{\TODO}[1]{%
    \par\noindent\textcolor{red}{\textbf{TODO:} #1}\par
}

\newcommand{\todobox}[1]{%
    \vspace{0.5em}
    \noindent
    \fcolorbox{red}{white}{%
        \parbox{0.96\linewidth}{\textcolor{red}{\textbf{TODO:} #1}}%
    }
    \vspace{0.5em}
}

\title{Hip Landmark Annotation Protocol}
\author{Orthopedic Driven Imaging}
\date{Version [X.Y] \\ Effective Date: [DATE]}

\begin{document}
\maketitle
\tableofcontents
\newpage

\section{Preliminary}

\todobox{This top-level section groups foundational protocol content before dataset 
construction and annotation operations begin. It is intended to contain Document Control, 
Purpose, Scope and Intended Use, and Governance, Change Control, and Sign-Off.}

\subsection{Document Control}

\subsubsection{Document Information}
\begin{itemize}[leftmargin=2em]
    \item Document title: Hip Landmark Annotation Protocol
    \item Document ID: [DOC-ID]
    \item Version: [X.Y]
    \item Effective date: [DATE]
    \item Supersedes: [OLDER VERSION IF APPLICABLE]
    \item Owner: [TEAM / FUNCTION]
\end{itemize}

\subsubsection{Version History}
\begin{longtable}{p{1.5cm} p{2.5cm} p{3cm} p{7cm}}
\toprule
Version & Date & Author(s) & Summary of changes \\
\midrule
0.1 & [DATE] & [NAME] & Initial protocol template \\
\bottomrule
\end{longtable}

\subsubsection{Approvals and Sign-Off}
\begin{itemize}[leftmargin=2em]
    \item Clinical lead: [NAME / TITLE]
    \item Annotation operations lead: [NAME / TITLE]
    \item ML / data science lead: [NAME / TITLE]
    \item Quality / regulatory lead: [NAME / TITLE]
\end{itemize}

\todobox{Complete the approval chain, signature requirements, release criteria, and 
document ownership. This addresses \textbf{Governance, change control, and sign-off} 
by making clear who is responsible for approving the protocol and who can authorize 
changes.}

\subsection{Purpose}
This document defines the protocol for creating and maintaining the reference standard 
for hip landmark annotation. It specifies the annotation scope, case eligibility rules, 
landmark definitions, annotator requirements, review and adjudication workflow, data 
quality expectations, and traceability requirements used to support safe, consistent, 
and reproducible labeling.

\todobox{Add a short paragraph stating that this document is the operational protocol 
for generating a fit-for-purpose reference standard for model development, validation, 
and ongoing lifecycle management. This helps satisfy \textbf{Scope and intended use} and 
\textbf{Governance} by stating what the document is for and how it will be used.}

\subsection{Scope and Intended Use}

\subsubsection{Protocol Scope}
This protocol applies to imaging studies used to generate landmark annotations, derived
 geometric constructs, and derived measurements for the hip landmarking tool.

The protocol covers:
\begin{itemize}[leftmargin=2em]
    \item primary landmarks,
    \item derived landmarks,
    \item measurements derived from primary and derived landmarks,
    \item review, adjudication, and QC procedures,
    \item dataset documentation and traceability.
\end{itemize}

It does not by itself define model architecture, training code, deployment criteria, or 
user interface behavior outside the annotation workflow.

\todobox{Insert the exact application \textbf{intended use statement}, including the 
clinical task, user population, care setting, and whether outputs are for planning, 
measurement assistance, decision support, post-op assessment, or research. Discuss how 
landmarks might be proxies for actual anatomical features (means to an end). This 
addresses \textbf{Scope and intended use} by tying the annotation protocol to the clinical 
purpose of the product.}

\subsubsection{Imaging Scope}
This protocol currently contemplates landmarks and measurements on AP and lateral imaging 
where applicable.

\todobox{Specify exactly which image types are in scope: e.g., AP radiographs, lateral hip 
radiographs, or other modalities. Also state whether this protocol applies to native hips, 
arthroplasty hips, revision cases, or all of the above. This addresses 
\textbf{Scope and intended use} and supports downstream 
\textbf{Data sources and case selection}.}

\subsubsection{Anatomical / Measurement Scope}
The current draft landmark set includes primary landmarks, derived landmarks, and 
measurements such as pelvic tilt, pelvic incidence, sacral slope, leg length, and 
femoral offset.

\todobox{Confirm which landmarks and measurements are in production scope versus 
exploratory / internal / future work. This helps address \textbf{Landmark dictionary} 
and ensures the protocol only governs the intended annotation outputs.}

\subsection{Governance, Change Control, and Sign-Off}

\subsubsection{Protocol Ownership}
This protocol is owned by [TEAM / FUNCTION].

\todobox{Define the document owner, operational owner, clinical owner, and quality owner,
 and state who is accountable for maintaining the protocol. This directly addresses 
 \textbf{Governance, change control, and sign-off}.}

\subsubsection{Change Control}
Changes to this protocol must be reviewed and approved before implementation.

\todobox{Define what constitutes a major versus minor change, who reviews each type of 
change, whether re-training or re-certification is required after changes, and how 
effective dates are managed. This directly addresses 
\textbf{Governance, change control, and sign-off}.}

\subsubsection{Sign-Off Rules}
No production annotation should proceed under an unapproved protocol version.

\todobox{Define the exact sign-off requirements and release gating rules before a 
protocol version can be used. This directly addresses 
\textbf{Governance, change control, and sign-off}.}

\section{Data}

\todobox{This top-level section groups dataset sourcing, eligibility, quality, and 
traceability requirements. It is intended to contain Data Sources and Case Selection, 
Inclusion and Exclusion Criteria, Data Quality Acceptance / Rejection Rules, and Dataset 
Traceability and Version Control.}

\subsection{Data Sources and Case Selection}

\subsubsection{Data Sources}
Cases used under this protocol may come from internal datasets, retrospective clinical 
archives, or external collaborators/providers, as permitted.

\todobox{List each allowed data source category and document where cases come from, what
 permissions govern their use, and whether each source is intended for development, 
 validation, or monitoring. This addresses \textbf{Data sources and case selection} 
 and supports \textbf{Dataset traceability and version control}.}

\subsubsection{Case Selection Strategy}
Cases should be selected according to the intended use population and target imaging 
conditions for the model.

\todobox{Add the exact case selection strategy: consecutive sampling, stratified sampling,
 enrichment of rare anatomy, balancing of implant/non-implant cases, balancing by site or
  projection, etc. State why the strategy was chosen and how it supports the intended use. 
  This addresses \textbf{Data sources and case selection} and 
  \textbf{Bias/representativeness plan}.}

\subsubsection{Site and Geography Strategy}
Cases should be sourced from sites and settings that are sufficiently representative of 
the intended use population and use environment.

\todobox{Specify the minimum number of collection sites, their geographic distribution, 
and why this configuration is fit for purpose. State whether data are sourced from 
distinct institutions / health systems, whether sites differ by care setting, equipment, 
and patient mix, and whether any site is reserved for external validation only. Define 
any cap on the proportion of cases contributed by a single site. This supports 
\textbf{Data sources and case selection} and \textbf{Bias/representativeness plan}.}

\subsubsection{Representativeness Targets}
Case accrual should be monitored against predefined representativeness targets relevant 
to the intended use population.

\todobox{Define which variables will be monitored for representativeness during accrual 
and dataset finalization, such as age, sex, geography, implant status, anatomy variation, 
disease severity, site, scanner/vendor, and projection/view. State any target ranges, 
minimum subgroup counts, or other acceptance criteria, and define how underrepresented 
subgroups will be remediated where feasible. This supports 
\textbf{Bias/representativeness plan} and \textbf{Data sources and case selection}.}

\subsubsection{Cohort Assignment and Independence}
Cases used for development, validation, testing, and post-deployment monitoring must be 
assigned to cohorts in a way that preserves independence and prevents information leakage.

\todobox{Define the cohort structure used under this protocol, including development, 
tuning, internal validation, locked test, and external validation sets where applicable. 
State the unit of assignment (e.g., patient, study, encounter, site), the rules that 
prevent the same patient or near-duplicate images from crossing cohorts, and any time- 
or site-based separation rules. This supports \textbf{Dataset traceability and version 
control} and \textbf{Bias/representativeness plan}.}

\subsubsection{Sequestration and Access Restrictions}
Access to evaluation datasets should be controlled to preserve the validity of downstream 
performance assessment.

\todobox{Specify which cohorts are sequestered, who is permitted to access them, what 
uses are prohibited, and how access is technically and operationally controlled. State 
whether model developers, annotators, adjudicators, or analysts may view locked test or 
external validation cases, and describe the approval and logging requirements for any 
exception. This supports \textbf{Dataset traceability and version control}, 
\textbf{Blinding and independence}, and \textbf{Governance}.}

\subsubsection{Case Metadata to Record}
At minimum, the following should be recorded where available:
\begin{itemize}[leftmargin=2em]
    \item study ID,
    \item series ID / image ID,
    \item projection / view,
    \item implant presence,
    \item annotator ID,
    \item protocol version,
    \item annotation timestamp.
\end{itemize}

\todobox{Expand this list to include all metadata required for stratification, 
subgroup analysis, and reproducibility, such as site, scanner/vendor, age bucket, 
sex, implant type, acquisition conditions, and cohort designation. This supports 
\textbf{Data sources and case selection}, \textbf{Bias/representativeness plan}, 
and \textbf{Dataset traceability and version control}.}

\subsection{Inclusion and Exclusion Criteria}

\subsubsection{Inclusion Criteria}
Eligible cases must contain anatomy and image characteristics sufficient for placement of 
the required landmarks for the applicable task.

\todobox{Add explicit inclusion criteria, such as required anatomy visible, accepted 
projections, minimum field-of-view, acceptable calibration marker presence, and whether 
bilateral or unilateral studies are acceptable. This directly addresses 
\textbf{Inclusion/exclusion criteria}.}

\subsubsection{Exclusion Criteria}
Cases may be excluded if anatomy is not sufficiently visible or if image quality or 
anatomical conditions make landmark placement invalid for the intended use.

\todobox{Add explicit exclusion criteria, such as severe motion artifact, truncation, 
extreme rotation, severe hardware obscuration, fracture patterns out of scope, post-op 
categories out of scope, fused anatomy, pediatric cases if excluded, or insufficient 
metadata. This directly addresses \textbf{Inclusion/exclusion criteria}.}

\subsubsection{Eligibility Review Process}
Eligibility should be assessed before production annotation begins.

\todobox{Add who performs eligibility review, whether it is done manually or 
algorithmically, what gets logged, and how borderline cases are escalated. This 
reinforces \textbf{Inclusion/exclusion criteria} and \textbf{Governance}.}

\subsection{Data Quality Acceptance / Rejection Rules}

\subsubsection{Image Quality Requirements}
Images used for annotation must be adequate for reliable placement of the required 
landmarks.

\todobox{Define objective acceptance and rejection rules for image quality, including 
exposure, contrast, motion, clipping/cropping, rotation/obliquity, artifacts, implant 
obscuration, etc. This directly addresses 
\textbf{Data quality acceptance/rejection rules}.}

\subsubsection{Anatomical Visibility Rules}
All landmarks required for a given task must either be visible enough to annotate or the 
case must be marked uncertain / not gradable according to this protocol.

\todobox{Add landmark-specific visibility thresholds or examples for what is considered 
sufficient anatomy visibility. This supports 
\textbf{Data quality acceptance/rejection rules} and 
\textbf{Uncertainty / not-visible / equivocal handling}.}

\subsubsection{Rejected Case Handling}
Rejected cases should not enter the finalized reference dataset without documented 
override.

\todobox{Specify how rejected cases are logged, whether they are retained in an 
excluded-case registry, who can override rejection, and how this decision is documented. 
This supports \textbf{Data quality acceptance/rejection rules}, \textbf{Traceability}, 
and \textbf{Governance}.}

\subsection{Dataset Traceability and Version Control}

\subsubsection{Traceability Requirements}
Each annotation event must be traceable to the source image, annotator, software version, 
and protocol version.

\todobox{Define the minimum traceability fields that must be stored for every case and 
every label edit, including image identifiers, annotator identity, timestamps, 
adjudication history, and finalization status. This directly addresses 
\textbf{Dataset traceability and version control}.}

\subsubsection{Version Control}
Datasets and protocol changes must be versioned.

\todobox{Define how protocol versions, landmark schema versions, dataset freezes, 
training/validation/test splits, and annotation revisions are version controlled. 
This directly addresses \textbf{Dataset traceability and version control}.}

\subsubsection{Change Logging}
Changes to annotations after initial completion must be auditable.

\todobox{Define who may edit finalized labels, how edits are logged, and when a changed 
label requires dataset re-release or retraining impact assessment. This supports 
\textbf{Dataset traceability and version control} and 
\textbf{Governance, change control, and sign-off}.}

\section{Annotation}

\todobox{This top-level section groups personnel requirements, annotation procedures, 
review workflows, uncertainty handling, QC, and bias / representativeness planning. 
It is intended to contain the core annotation operations content.}

\subsection{Annotator Qualifications}

\subsubsection{Annotator Roles}
The annotation program may include:
\begin{itemize}[leftmargin=2em]
    \item primary annotators,
    \item secondary reviewers,
    \item adjudicators,
    \item clinical oversight reviewers.
\end{itemize}

\todobox{Define minimum qualifications for each role, including clinical background, 
anatomy knowledge, prior annotation experience, and any role-specific credentialing 
requirements. This directly addresses \textbf{Annotator qualifications}.}

\subsubsection{Qualification Requirements}
Annotators must be qualified for the tasks they perform.

\todobox{State the minimum education, training, and competency requirements for each role. 
Also define when non-clinicians may annotate independently versus when clinical 
supervision is required. This directly addresses \textbf{Annotator qualifications}.}

\subsection{Annotator Training and Certification}

\subsubsection{Initial Training}
All annotators must complete protocol-specific training before contributing production 
labels.

\todobox{Add the initial training package contents: anatomy review, landmark definitions,
 software usage, common edge cases, escalation rules, and practice set review. 
 This directly addresses \textbf{Annotator training/certification}.}

\subsubsection{Calibration}
Annotators should complete a calibration round on a predefined set of cases.

\todobox{Define the calibration dataset, required agreement thresholds, and remediation 
steps if the annotator does not meet threshold. This directly addresses 
\textbf{Annotator training/certification} and supports 
\textbf{QC and reliability assessment}.}

\subsubsection{Certification and Re-Certification}
Only certified (trained) annotators may produce production annotations.

\todobox{Define the pass/fail criteria for certification, re-certification frequency, 
retraining triggers, and documentation requirements. This directly addresses 
\textbf{Annotator training/certification}.}

\subsection{Annotation Environment and General Rules}

\subsubsection{Annotation Software and Display Environment}
Annotation is performed in the approved annotation environment.

\todobox{Specify software/tool version, controls, software tool usage, preprocessing 
rules, etc. Include a manual for how to use the annotation tool. This supports 
\textbf{Annotation instructions by landmark} and 
\textbf{Dataset traceability and version control}.}

\subsubsection{General Annotation Principles}
\begin{enumerate}[leftmargin=2em]
    \item Annotate only on approved image types.
    \item Do not infer anatomy not visible on the image.
    \item Follow landmark-specific definitions exactly.
    \item Escalate uncertain cases rather than guessing.
\end{enumerate}

\todobox{Expand these into the universal annotation rules used across all landmarks, 
including rules about interpolation, uncertainty, common judgments/scenarios, etc. 
This supports \textbf{Annotation instructions by landmark}.}

\subsection{Landmark Dictionary}

This section defines the canonical landmark inventory governed by the protocol. 
Each landmark entry should state the acronym, underlying anatomical feature(s), view(s), 
prediction geometry, whether it is primary or derived, and, if primary, unambiguous 
placement rules, edge-case handling, and when to escalate. Note that if a landmark has 
dual points, lines, or other features corresponding to the left and right sides of the 
anatomy, then the acronym used to refer to that particular geometric feature will be the 
normal acronym name prepended with an L or R, as appropriate. 

\todobox{Populate this section with the finalized inventory of all primary and derived 
landmarks and measurements, and clearly distinguish production, internal-only, 
exploratory, and deferred items. This directly addresses \textbf{Landmark dictionary}.}

\subsubsection{Primary Landmarks}
\todobox{Make this complete.}
\begin{itemize}[leftmargin=2em]
    \item Acetabular Center
    \item Anterior Superior Iliac Spine
    \item Femoral Head Center
    \item Ischial Tuberositiy
    \item Superior Greater Trochanter
    \item Lateral Greater Trochanter
    \item Lesser Trochanter
    \item Pelvic Outlet Diameter
    \item Superior Margin of the Pubic Symphysis
    \item Pelvic Teardrop
    \item Femoral Axis

\end{itemize}

\subsubsection{Derived Landmarks}
\todobox{Are these correct? Make this complete.}
\begin{itemize}[leftmargin=2em]
    \item Transischial Line
    \item Trans-SI Line
    \item Sacral Endplate Line
    \item Lesser Trochanter Line
    \item Distal Teardrop Line
    \item Sacral Endplate -- Femoral Head Center Line
    \item Anteinclination Line
    \item Femoral Anatomical Axis
\end{itemize}

\subsubsection{Measurements}
\todobox{Are these correct? Make this complete.}
\begin{itemize}[leftmargin=2em]
    \item Sagittal Pelvic Tilt
    \item Pelvic Incidence
    \item Sacral Slope
    \item Sacral Apex Angle
    \item Anteinclination Angle
    \item Sacroacetabular Angle
    \item Pelvic Acetabular Angle
    \item Pelvic Femoral Angle
    \item Leg Length
    \item Femoral Diaphyseal Length
    \item Femoral Offset
\end{itemize}

\subsection{Annotation Instructions by Landmark}

For each landmark, this protocol should define:
\begin{itemize}[leftmargin=2em]
    \item exact anatomical definition,
    \item applicable view(s),
    \item point placement rule,
    \item tie-break rules,
    \item edge cases,
    \item non-examples,
    \item escalation triggers,
    \item whether user modification is allowed.
\end{itemize}

\todobox{Create a dedicated subsection for each landmark and write the full operational 
instruction for annotators. Include exact point placement rules, edge-case handling, 
and examples of correct vs incorrect placement. This directly addresses 
\textbf{Annotation instructions by landmark}.}

\subsubsection{Example Landmark Entry Template}
\noindent
\textbf{Name:} [Landmark Name]

\noindent
\textbf{Acronynm:} [LN]

\noindent
\textbf{Underlying Anatomical Feature(s):} [Anatomical definition]

\noindent
\textbf{View(s):} [AP / lateral / etc.]

\noindent
\textbf{Prediction Geometry:} [N points/ lines / etc.]

\noindent
\textbf{Type:} [Primary/derived]

\noindent
\textbf{Placement Rules:} [Detailed instruction]

\noindent
\textbf{Edge-Case Handling:} [Detailed instruction]

\noindent
\textbf{Escalate When:} [Detailed instruction]


% Acetabular Center
\subsubsection{Acetabular Center}
\noindent
\textbf{Name:} Acetabular Center

\noindent
\textbf{Acronynm:} AC

\noindent
\textbf{Underlying Anatomical Feature(s):} [Anatomical definition]

\noindent
\textbf{View(s):} AP

\noindent
\textbf{Prediction Geometry:} 2 points (left and right)

\noindent
\textbf{Type:} Primary

\noindent
\textbf{Placement Rules:} [Detailed instruction]

\noindent
\textbf{Edge-Case Handling:} [Detailed instruction]

\noindent
\textbf{Escalate When:} [Detailed instruction]

% Anterior Superior Iliac Spine
\subsubsection{Anterior Superior Iliac Spine}
\noindent
\textbf{Name:} Anterior Superior Iliac Spine

\noindent
\textbf{Acronynm:} ASIS

\noindent
\textbf{Underlying Anatomical Feature(s):} [Anatomical definition]

\noindent
\textbf{View(s):} AP

\noindent
\textbf{Prediction Geometry:} 2 points (left and right)

\noindent
\textbf{Type:} Primary

\noindent
\textbf{Placement Rules:} [Detailed instruction]

\noindent
\textbf{Edge-Case Handling:} [Detailed instruction]

\noindent
\textbf{Escalate When:} [Detailed instruction]

% Femoral Head Center
\subsubsection{Femoral Head Centerr}
\noindent
\textbf{Name:} Femoral Head Centerr

\noindent
\textbf{Acronynm:} FHC

\noindent
\textbf{Underlying Anatomical Feature(s):} [Anatomical definition]

\noindent
\textbf{View(s):} AP

\noindent
\textbf{Prediction Geometry:} 2 points (left and right)

\noindent
\textbf{Type:} Primary

\noindent
\textbf{Placement Rules:} [Detailed instruction]

\noindent
\textbf{Edge-Case Handling:} [Detailed instruction]

\noindent
\textbf{Escalate When:} [Detailed instruction]

% Ischial Tuberositiy
\subsubsection{Ischial Tuberositiy}
\noindent
\textbf{Name:} Ischial Tuberositiy

\noindent
\textbf{Acronynm:} IT

\noindent
\textbf{Underlying Anatomical Feature(s):} [Anatomical definition]

\noindent
\textbf{View(s):} AP

\noindent
\textbf{Prediction Geometry:} 2 points (left and right)

\noindent
\textbf{Type:} Primary

\noindent
\textbf{Placement Rules:} [Detailed instruction]

\noindent
\textbf{Edge-Case Handling:} [Detailed instruction]

\noindent
\textbf{Escalate When:} [Detailed instruction]

% Superior Greater Trochanter
\subsubsection{Superior Greater Trochanter}
\noindent
\textbf{Name:} Superior Greater Trochanter

\noindent
\textbf{Acronynm:} SGT

\noindent
\textbf{Underlying Anatomical Feature(s):} [Anatomical definition]

\noindent
\textbf{View(s):} AP

\noindent
\textbf{Prediction Geometry:} 2 points (left and right)

\noindent
\textbf{Type:} Primary

\noindent
\textbf{Placement Rules:} [Detailed instruction]

\noindent
\textbf{Edge-Case Handling:} [Detailed instruction]

\noindent
\textbf{Escalate When:} [Detailed instruction]

% Lateral Greater Trochanter
\subsubsection{Lateral Greater Trochanter}
\noindent
\textbf{Name:} Lateral Greater Trochanter

\noindent
\textbf{Acronynm:} LGT

\noindent
\textbf{Underlying Anatomical Feature(s):} [Anatomical definition]

\noindent
\textbf{View(s):} AP

\noindent
\textbf{Prediction Geometry:} 2 points (left and right)

\noindent
\textbf{Type:} Primary

\noindent
\textbf{Placement Rules:} [Detailed instruction]

\noindent
\textbf{Edge-Case Handling:} [Detailed instruction]

\noindent
\textbf{Escalate When:} [Detailed instruction]

% Lesser Trochanter
\subsubsection{Lesser Trochanter}
\noindent
\textbf{Name:} Lesser Trochanter

\noindent
\textbf{Acronynm:} LT

\noindent
\textbf{Underlying Anatomical Feature(s):} [Anatomical definition]

\noindent
\textbf{View(s):} AP

\noindent
\textbf{Prediction Geometry:} 2 points (left and right)

\noindent
\textbf{Type:} Primary

\noindent
\textbf{Placement Rules:} [Detailed instruction]

\noindent
\textbf{Edge-Case Handling:} [Detailed instruction]

\noindent
\textbf{Escalate When:} [Detailed instruction]

% Pelvic Outlet Diameter
\subsubsection{Pelvic Outlet Diameter}
\noindent
\textbf{Name:} Pelvic Outlet Diameter

\noindent
\textbf{Acronynm:} POD

\noindent
\textbf{Underlying Anatomical Feature(s):} [Anatomical definition]

\noindent
\textbf{View(s):} AP

\noindent
\textbf{Prediction Geometry:} 2 points (left and right)

\noindent
\textbf{Type:} Primary

\noindent
\textbf{Placement Rules:} [Detailed instruction]

\noindent
\textbf{Edge-Case Handling:} [Detailed instruction]

\noindent
\textbf{Escalate When:} [Detailed instruction]

% Superior Margin of the Pubic Symphysis
\subsubsection{Superior Margin of the Pubic Symphysis}
\noindent
\textbf{Name:} Superior Margin of the Pubic Symphysis

\noindent
\textbf{Acronynm:} SMPS

\noindent
\textbf{Underlying Anatomical Feature(s):} [Anatomical definition]

\noindent
\textbf{View(s):} AP

\noindent
\textbf{Prediction Geometry:} 1 point

\noindent
\textbf{Type:} Primary

\noindent
\textbf{Placement Rules:} [Detailed instruction]

\noindent
\textbf{Edge-Case Handling:} [Detailed instruction]

\noindent
\textbf{Escalate When:} [Detailed instruction]

% Pelvic Teardrop
\subsubsection{Pelvic Teardrop}
\noindent
\textbf{Name:} Pelvic Teardrop

\noindent
\textbf{Acronynm:} PT

\noindent
\textbf{Underlying Anatomical Feature(s):} [Anatomical definition]

\noindent
\textbf{View(s):} AP

\noindent
\textbf{Prediction Geometry:} 2 points (left and right)

\noindent
\textbf{Type:} Primary

\noindent
\textbf{Placement Rules:} [Detailed instruction]

\noindent
\textbf{Edge-Case Handling:} [Detailed instruction]

\noindent
\textbf{Escalate When:} [Detailed instruction]

% Femoral Axis
\subsubsection{Femoral Axis}
\noindent
\textbf{Name:} Femoral Axis

\noindent
\textbf{Acronynm:} FA

\noindent
\textbf{Underlying Anatomical Feature(s):} [Anatomical definition]

\noindent
\textbf{View(s):} AP

\noindent
\textbf{Prediction Geometry:} 2 lines (left and right)

\noindent
\textbf{Type:} Primary

\noindent
\textbf{Placement Rules:} [Detailed instruction]

\noindent
\textbf{Edge-Case Handling:} [Detailed instruction]

\noindent
\textbf{Escalate When:} [Detailed instruction]

\todobox{Repeat the above template for every landmark in scope. This ensures the protocol 
is operationally complete and satisfies \textbf{Annotation instructions by landmark}.}

\subsection{Blinding and Independence}

\subsubsection{Blinding Policy}
Annotations should be collected in a way that minimizes avoidable bias.

\todobox{Specify whether annotators are blinded to model predictions, prior annotations, 
clinical labels, implant metadata, or each other's work. State what information is 
available during annotation and why. This directly addresses 
\textbf{Blinding and independence}.}

\subsubsection{Independence of Reads}
When multiple reads are obtained, they should be independent unless the protocol 
explicitly calls for consensus annotation.

\todobox{Define when reads must be independent, when reviewers may see prior annotations,
 and how independence is maintained operationally. This directly addresses 
 \textbf{Blinding and independence}.}

\subsection{Adjudication and Consensus Rules}

\subsubsection{Triggers for Review}
Cases may be reviewed or adjudicated when disagreement, uncertainty, or QC failure occurs.

\todobox{Define exactly what triggers adjudication: disagreement thresholds, missing 
landmarks, difficult anatomy, low confidence, outlier measurements, or random QC sampling. 
This directly addresses \textbf{Adjudication and consensus rules}.}

\subsubsection{Adjudication Workflow}
A designated reviewer or panel should resolve flagged cases.

\todobox{Describe the stepwise adjudication workflow: who reviews first, who makes the 
final decision, whether consensus or single-expert override is used, and how final labels 
are chosen. This directly addresses \textbf{Adjudication and consensus rules}.}

\subsubsection{Consensus Rules}
Consensus rules govern how multiple annotations are resolved into the final reference 
standard.

\todobox{Specify whether final labels are produced by majority vote, adjudicator 
overwrite, geometric averaging, or another method, and justify why this approach is 
fit for purpose. This directly addresses \textbf{Adjudication and consensus rules}.}

\subsection{Uncertainty, Not-Visible, and Equivocal Handling}

\subsubsection{Not Visible}
A landmark may be marked not visible when the anatomy needed for valid placement is not 
adequately visualized.

\todobox{Define exactly when a landmark can be marked not visible, whether the case 
remains usable, and whether alternative landmarks or fallback measurements are allowed. 
This directly addresses \textbf{Uncertainty / not-visible / equivocal handling}.}

\subsubsection{Equivocal Cases}
Some cases may have ambiguous anatomy or multiple plausible placements.

\todobox{Define how equivocal cases are identified, labeled, escalated, and resolved. 
State whether annotators may place a best estimate or must defer to adjudication. This 
directly addresses \textbf{Uncertainty / not-visible / equivocal handling}.}

\subsubsection{Confidence / Uncertainty Capture}
If uncertainty is recorded, it should be captured in a structured way.

\todobox{Decide whether to collect confidence scores, uncertainty tags, or reason codes, 
and document how these are used in downstream QC or model development. This directly 
addresses \textbf{Uncertainty / not-visible / equivocal handling}.}

\subsection{QC and Reliability Assessment}

\subsubsection{Routine QC}
Routine QC should confirm that annotations are complete, plausible, and consistent with 
protocol rules.

\todobox{Define the QC plan, including completeness checks, landmark plausibility checks, 
left-right consistency checks, derived-measurement sanity checks, and manual spot review 
rates. This directly addresses \textbf{QC and reliability assessment}.}

\subsubsection{Reliability Assessment}
Reliability should be evaluated periodically and documented.

\todobox{Define inter-annotator and intra-annotator reliability methods, the sample size 
for repeated reads, the metrics used (e.g., point distance error, ICC), and acceptable 
thresholds. This directly addresses \textbf{QC and reliability assessment}.}

\subsubsection{Corrective Action}
If QC or reliability thresholds are not met, corrective action is required.

\todobox{State what happens when QC fails: retraining, reannotation, reviewer escalation, 
temporary suspension of annotator privileges, or protocol revision. This supports 
\textbf{QC and reliability assessment} and \textbf{Governance}.}

\subsection{Bias and Representativeness Plan}

\subsubsection{Target Population Representation}
The dataset should represent the intended use population and key clinically relevant 
subgroups.

\todobox{Define the subgroup variables that matter for your use case, such as age, sex, 
implant status, anatomy variation, site, scanner/vendor, disease severity, and projection 
differences. Explain how the case selection process will ensure adequate representation. 
This directly addresses \textbf{Bias/representativeness plan}.}

\subsubsection{Monitoring for Imbalance}
Dataset composition should be monitored during case accrual and dataset finalization.

\todobox{Add a plan for tracking subgroup counts, identifying underrepresented 
populations, and correcting imbalance where feasible. This directly addresses 
\textbf{Bias/representativeness plan}.}

\subsubsection{Performance-Relevant Subgroups}
Where relevant, annotations and derived measurements should support subgroup-level 
evaluation.

\todobox{Add a section identifying which subgroups will later be used in model 
performance evaluation and why they matter clinically. This directly addresses 
\textbf{Bias/representativeness plan}.}

\end{document}